\chapter{Infinitive Constructions}

An infinitive can appear without \textit{zu} or as part of a construction with
\textit{zu}. The position of \textit{zu}, the understood subject, and the
comma are the most important points for learners.

\section{The infinitive without \textit{zu}}

No \textit{zu} is used after modal verbs or \textit{werden}, nor in several
common constructions with verbs of perception, movement, and causation.

\begin{center}
\small
\rowcolors{2}{referenceBlueBg}{white}
\begin{tabularx}{\textwidth}{@{}>{\bfseries}l X@{}}
\toprule
Construction & Example \\
\midrule
Modal verb & \textit{Ich muss heute arbeiten.} \\
Futur with \textit{werden} & \textit{Wir werden morgen kommen.} \\
\textit{lassen} & \textit{Lass mich bitte ausreden.} \\
Perception & \textit{Ich höre die Kinder spielen.} \\
Movement + activity & \textit{Wir gehen heute schwimmen.} \\
\textit{lernen} & \textit{Das Kind lernt lesen.} \\
\bottomrule
\end{tabularx}
\end{center}

\section{The infinitive with \textit{zu}}

Many verbs, adjectives, nouns, and impersonal expressions introduce a
\textit{zu}-infinitive.

\begin{center}
\small
\rowcolors{2}{referenceGreenBg}{white}
\begin{tabularx}{\textwidth}{@{}>{\bfseries}l X@{}}
\toprule
Trigger & Example \\
\midrule
Verb & \textit{Ich versuche, jeden Tag Deutsch zu lernen.} \\
Adjective & \textit{Es ist wichtig, genug zu schlafen.} \\
Noun & \textit{Ich habe den Wunsch, in Berlin zu studieren.} \\
Expression with anticipatory \textit{es}
  & \textit{Ich finde es schwer, früh aufzustehen.} \\
\bottomrule
\end{tabularx}
\end{center}

The understood subject of the infinitive is usually supplied by the main
clause. Compare:

\begin{itemize}[leftmargin=*]
  \item \textit{Ich hoffe, bald zu reisen.} --- \textit{ich} will travel.
  \item \textit{Ich bitte Paul, bald zu kommen.} --- \textit{Paul} should come.
\end{itemize}

If the two actions have different explicit subjects, use a finite subordinate
clause instead: \textit{Ich hoffe, dass Paul bald kommt.}

\section{Where \textit{zu} goes}

\begin{center}
\small
\rowcolors{2}{referenceYellowBg}{white}
\begin{tabularx}{\textwidth}{@{}>{\bfseries}l l l X@{}}
\toprule
Verb type & Infinitive & With \textit{zu} & Example \\
\midrule
Simple & lernen & zu lernen & \textit{Ich versuche, Deutsch zu lernen.} \\
Inseparable & verstehen & zu verstehen & \textit{Es ist schwer, das zu verstehen.} \\
Verb in \textit{-ieren} & reservieren & zu reservieren & \textit{Ich vergaß, zu reservieren.} \\
Separable & anrufen & anzurufen & \textit{Ich verspreche, dich anzurufen.} \\
Separable & aufstehen & aufzustehen & \textit{Ich versuche, früh aufzustehen.} \\
\bottomrule
\end{tabularx}
\end{center}

With a separable verb, \textit{zu} is inserted between the prefix and the verb
stem: \textit{an + zu + rufen = anzurufen}.

\section{Purpose, absence, and alternatives}

\begin{center}
\small
\rowcolors{2}{referenceRedBg}{white}
\begin{tabularx}{\textwidth}{@{}>{\bfseries}l l X@{}}
\toprule
Construction & Meaning & Example \\
\midrule
um \ldots{} zu & in order to & \textit{Ich lerne viel, um die Prüfung zu bestehen.} \\
ohne \ldots{} zu & without doing & \textit{Er ging, ohne sich zu verabschieden.} \\
statt/anstatt \ldots{} zu & instead of doing
  & \textit{Sie fährt Rad, statt das Auto zu nehmen.} \\
\bottomrule
\end{tabularx}
\end{center}

Use \textit{um \ldots{} zu} when the understood subject is the same in both
parts. Use \textit{damit} plus a finite clause when the subjects differ:

\begin{center}
\textit{Ich spreche langsam, um verständlich zu sein.}

\textit{Ich spreche langsam, damit die Kinder mich verstehen.}
\end{center}

\section{\textit{brauchen + zu}}

In standard usage, \textit{brauchen + zu} normally occurs with a negative or
limiting expression and has a modal meaning:

\begin{center}
\textit{Du brauchst heute nicht zu kommen.}
\qquad
\textit{Du brauchst nur anzurufen.}
\end{center}

\section{Perfect infinitive}

To show that the infinitive action happened earlier, use Partizip II plus
\textit{zu haben} or \textit{zu sein}. The auxiliary follows the same choice
as in the Perfekt.

\begin{itemize}[leftmargin=*]
  \item \textit{Sie behauptet, den Brief geschrieben \textbf{zu haben}.}
  \item \textit{Er scheint früh angekommen \textbf{zu sein}.}
\end{itemize}

\section{Commas with infinitive groups}

Under the current official spelling rules, an expanded, clause-like
infinitive group is separated by a comma:
\textit{Sie verspricht, morgen pünktlich zu kommen.} A comma is also required
with \textit{um, ohne, statt, anstatt, außer}, and \textit{als}, and when the
group is inserted into a sentence it is enclosed by a pair of commas.

A bare, unexpanded \textit{zu}-infinitive can be written without a comma:
\textit{Sie verspricht zu kommen.} No comma separates an integrated infinitive
construction with verbs such as \textit{brauchen} or \textit{scheinen}:
\textit{Du brauchst nicht zu warten}; \textit{Er scheint zu schlafen}.

\begin{tcolorbox}[
  colback=referenceBlueBg,
  colframe=genderMasculine,
  title={Quick choice},
  fonttitle=\bfseries
]
\begin{itemize}[leftmargin=*,nosep]
  \item After a modal verb: infinitive without \textit{zu}.
  \item Same subject and a purpose: \textit{um \ldots{} zu}.
  \item Different subjects: usually \textit{damit} or another finite clause.
  \item Separable verb: put \textit{zu} between prefix and stem.
  \item Expanded clause-like infinitive group: separate it with a comma.
\end{itemize}
\end{tcolorbox}
